\chapter{Instrucciones de uso de la plantilla}

\epigraph{La perfección no se alcanza cuando no hay nada más que añadir, sino cuando no hay nada más que quitar.}{Antoine de Saint-Exupéry}

\lettrine[lraise=-0.1, lines=2, loversize=0.2]{E}{sta} nueva plantilla nace con el fin de ofrecer una versión simplificada del formato de publicaciones de la ETSi tratando de respetar en la mayor medida posible su estética y usabilidad, pero priorizando la velocidad de compilación para poder utilizarse con editores web como Overleaf, tan populares en la actualidad, pero que presentan limitaciones de cómputo significativas.

Este capítulo presenta sus posibles usos. De momento, la plantilla sólo está disponible para Trabajos Fin de Grado o Máster. 
A fecha de redacción de este documento, está preparada para ser compilada con \textit{pdflatex}, y ha sido probada en TeX Live 2025 tanto en local como en Overleaf.

\section{Estructura básica del documento}

Los documentos en \LaTeX{} presentan dos partes diferenciadas: 
\begin{itemize}\itemsep 0em
    \item El \textbf{preámbulo} le permite definir la clase del documento (en este caso, \textit{etsibook}, los paquetes a utilizar, y la personalización de la portada, entre otros. Abarca todo el contenido hasta la sentencia \lstinline|\begin{document}|.
    \item El \textbf{cuerpo} incluye los contenidos de su documento encerrados entre \lstinline|\begin{document}| y \lstinline|\end{document}|.
\end{itemize}

En el código fuente de esta plantilla encontrará comentarios, precedidos por el carácter \lstinline|%|, que le guiarán en la personalización del documento. Adicionalmente, a continuación se presenta brevemente las diferentes posibilidades de la plantilla.

\subsection{El preámbulo}

El preámbulo comienza con la instrucción \lstinline|\|\lstinline|documentclass{etsibook}[opciones]|, que le permite personalizar algunas características básicas del documento: tamaño de la fuente, geometría del documento, etc. Para los trabajos Fin de Grado o Máster se recomienda mantener los valores por defecto (fuente a 10pt, y formato A4).

A continuación, el usuario puede incluir los paquetes que necesite para la elaboración del documento con la sentencia \lstinline|\usepackage{paquete}[opciones]|. Es recomendable incluir sólo aquellos paquetes imprescindibles para mejorar los tiempos de compilación del documento. Por ejemplo, si no necesita incluir fragmentos de código, comente la importación del paquete \textit{listings} anteponiendo el carácter \lstinline|%| a la línea \lstinline|\usepackage{listings}|. 

El siguiente bloque del preámbulo le permite definir la portada del documento y sus metadatos. Para ello, la plantilla ofrece los siguientes comandos:
\begin{lstlisting}[language=tex, caption={Bloque de personalización de la portada}]
\header{Trabajo Fin de Grado\\%
  Grado en Ingeniería XXX}
\title{Escuela Técnica Superior de Ingeniería: Formato de Publicación para Overleaf}
\author{Nombre Apellidos}
\authorlabel{Autor}
\director{Tutor}{Nombre Apellidos}{Categoría} 
\date{Sevilla, 2025}
\affiliation{Dpto. de Ingeniería XXX\\%
        Escuela Técnica Superior de Ingeniería\\%
        Universidad de Sevilla}
\optlogo{etsi/logo-dpto-placeholder.pdf} 
\end{lstlisting}

En general basta con cambiar estos valores por defecto para adecuarlos a su documento: titulación, título del documento, nombre y apellidos del autor y tutores. Algunas consideraciones son:
\begin{itemize}\itemsep 0em
    \item La instrucción \lstinline|\director| puede incluirse varias veces si tiene varios directores.
    \item La orden \lstinline|\optlogo| es opcional y le permite incluir un segundo logotipo en la portada. Puede comentar esta línea si no la utiliza.
    \item Puede incluir saltos de línea con \lstinline|\\| si lo necesita.
    \item Si uno de los campos presenta un texto muy largo, puede anteponer \lstinline|\smaller{}| a su texto para reducir ligeramente el tamaño de la fuente.
\end{itemize}

Finalmente puede personalizar otros comandos o paquetes. Por ejemplo, esta plantilla prepara el glosario utilizando el paquete \lstinline|glossaries|.

\subsection{El cuerpo}

El cuerpo del documento incluye sus contenidos entre las órdenes \lstinline|\begin{document}| y \lstinline|\end{document}|. En él, se distinguen cuatro bloques principales: 
\begin{itemize}\itemsep 0em
    \item La \textbf{cubierta}, que se construye con la orden \lstinline{\maketitle} y debe ser lo primero en el cuerpo del documento.
    \item El \textbf{contenido preliminar} (agradecimientos, resumen, \textit{abstract}, etc.), que se debe incluir tras la orden \lstinline{\frontmatter}, ya que este comando establece una numeración especial. En este bloque, los capítulos deben comenzarse con la orden \lstinline{\addchap} para que no se numeren pero aparezcan en el índice.
    \item El \textbf{contenido principal} del documento, que debe incluirse tras la orden \lstinline{\mainmatter} que vuelve a personalizar la numeración, el formato de la cabecera y del pie de página. En este bloque, los capítulos deben comenzarse con la orden \lstinline{\chapter}.
    \item El \textbf{contenido final} (índices de figuras, anexos, etc.) debe ubicarse tras la orden \lstinline{\backmatter}.
\end{itemize}

\section{Catálogo de estilos}

Esta sección muestra los principales elementos tipográficos disponibles de la plantilla, a modo de catálogo visual y didáctico. De este modo, puede comparar el código \LaTeX{} usado para reproducir el elemento. 

\subsection{Formato de texto}

El texto ordinario se escribe directamente en el cuerpo del documento, sin ningún comando especial. Para aplicar estilos básicos, \LaTeX{} dispone de los siguientes comandos de la Tabla \ref{tab:formato}.

\begin{table}[htbp]
\caption{Comandos de formato de texto}\label{tab:formato}
\vspace{0.1cm}
\centering
\begin{tabular}{lll}
\hline
\textbf{Comando} & \textbf{Resultado} & \textbf{Uso} \\
\hline
\lstinline|\textbf{}| & \textbf{negrita} & Énfasis fuerte \\
\lstinline|\textit{}| & \textit{cursiva} & Resalto, términos extranjeros \\
\lstinline|\texttt{}| & \texttt{monoespaciado} & Código, rutas de fichero \\
\lstinline|\emph{}|   & \emph{énfasis}   & Énfasis contextual \\
\lstinline|\underline{}| & \underline{subrayado} & Usar con moderación \\[0.05cm] 
\hline
\end{tabular}
\end{table}

\subsubsection*{Tamaños de fuente}

La plantilla define el tamaño de fuente base en las opciones de la clase. Es posible variar el tamaño localmente con los siguientes comandos (de menor a mayor):

\begin{center}
{\tiny tiny} \quad {\scriptsize scriptsize} \quad {\footnotesize footnotesize} \quad {\small small} \quad {\normalsize normalsize} 

{\large large} \quad {\Large Large} \quad {\LARGE LARGE} 

{\huge huge} \quad {\Huge Huge}
\end{center}

Además, la plantilla incluye los comandos relativos \lstinline|\larger{}| y \lstinline|\smaller{}| para ajustes finos en la portada, aunque también pueden usarse en el texto si fuera necesario.

\subsection{Estructura del texto}

\subsubsection*{Epígrafe}

Puede iniciar un capítulo con un epígrafe usando el comando \lstinline|\epigraph{}{}|, tal como se muestra al comienzo de este capítulo.

\subsubsection*{Letra capital (\textit{lettrine})}

\lettrine[lraise=-0.1, lines=2, loversize=0.2]{E}{l} inicio de un capítulo puede destacarse con una letra capital decorativa usando el paquete \textit{lettrine}. Los parámetros de \lstinline{\lettrine} permiten configurar el número de líneas y el tamaño de la letra, entre otros.

\subsubsection*{Notas al pie}

Las notas al pie de página se insertan de forma sencilla con el comando \lstinline|\footnote{}|, directamente en el punto del texto donde se quiere insertar la llamada\footnote{Esta es una nota al pie de ejemplo.}.

\subsection{Listas}

Se han personalizado los siguientes tipos de lista. Las no numeradas, mediante el entorno \lstinline{itemize}:
\begin{itemize}\itemsep 0em
    \item Primer elemento de la lista.
    \item Segundo elemento de la lista.
    \item Tercer elemento, con un subnivel:
    \begin{itemize}
        \item Primer subelemento.
        \item Segundo subelemento.
    \end{itemize}
\end{itemize}

Y las numeradas, usando el entorno \lstinline{enumerate}:
\begin{enumerate}\itemsep 0em
    \item Primer paso del procedimiento.
    \item Segundo paso del procedimiento.
    \item Tercer paso, que a su vez se divide en:
    \begin{enumerate}
        \item Primera sub-tarea.
        \item Segunda sub-tarea.
    \end{enumerate}
\end{enumerate}

\subsection{Ecuaciones}

Las expresiones matemáticas se escriben con el paquete \textit{amsmath}. Puede incluir matemáticas en línea encerrando la expresión entre signos de dólar, por ejemplo: la transformada de Fourier es $F(\omega) = \int_{-\infty}^{\infty} f(t)\,e^{-j\omega t}\,dt$.

\subsubsection*{Ecuación numerada}

El entorno \textit{equation} produce una ecuación numerada y centrada:

\begin{equation}
    E = mc^{2}
    \label{eq:einstein}
\end{equation}

\subsubsection*{Sistema de ecuaciones (\textit{align})}

El entorno \textit{align} permite alinear varias ecuaciones. Las ecuaciones que no deban numerarse llevan \lstinline|\nonumber|:

\begin{align}
    \nabla \cdot \mathbf{E} &= \frac{\rho}{\varepsilon_0} \label{eq:gauss} \\
    \nabla \cdot \mathbf{B} &= 0 \nonumber \\
    \nabla \times \mathbf{E} &= -\frac{\partial \mathbf{B}}{\partial t} \nonumber \\
    \nabla \times \mathbf{B} &= \mu_0\mathbf{J} + \mu_0\varepsilon_0\frac{\partial \mathbf{E}}{\partial t} \label{eq:ampere}
\end{align}

\subsubsection*{Matrices}

Las matrices se representan con los entornos \textit{pmatrix} (paréntesis), \textit{bmatrix} (corchetes) o \textit{vmatrix} (barras). Se insertan dentro de un entorno matemático:

\begin{equation}
    \mathbf{A} = \begin{bmatrix}
        a_{11} & a_{12} & a_{13} \\
        a_{21} & a_{22} & a_{23} \\
        a_{31} & a_{32} & a_{33}
    \end{bmatrix}
    \label{eq:matriz}
\end{equation}

\subsection{Tablas}

Las tablas se definen con el entorno flotante \textit{table} (que gestiona la posición) y el entorno \textit{tabular} (que define el contenido). Note que en esta plantilla el \lstinline|\caption{}| va antes del contenido, lo que es la convención habitual para tablas.

\begin{table}[htbp]
\caption{Comparativa de protocolos de transporte}\label{tab:protocolos}
\vspace{0.1cm}
\centering
\begin{tabular}{llll}
\hline
\textbf{Característica} & \textbf{TCP} & \textbf{UDP} & \textbf{SCTP} \\
\hline
Orientado a conexión   & Sí  & No  & Sí  \\
Entrega garantizada    & Sí  & No  & Sí  \\
Control de flujo       & Sí  & No  & Sí  \\
Multihoming            & No  & No  & Sí  \\
\hline
\end{tabular}
\end{table}

\subsection{Figuras}

Las figuras son elementos flotantes definidos con el entorno \textit{figure}. El \lstinline|\caption{}| se sitúa \textbf{después} de la imagen. El parámetro de posición \lstinline{[htbp]} indica a \LaTeX{} que intente colocar la figura aquí (\textit{here}), al principio de página (\textit{top}), al final de página (\textit{bottom}) o en una página dedicada (\textit{page}).

\begin{figure}[htbp]
    \centering
    \includegraphics[width=0.15\linewidth]{etsi/logo-etsi.pdf}
    \caption{Logotipo de la Escuela Técnica Superior de Ingeniería}
    \label{fig:logoetsi-cat}
\end{figure}

\subsection{Fragmentos de código}

El paquete \textit{listings} permite incluir fragmentos de código con resaltado sintáctico. Puede especificar el lenguaje con la opción \lstinline{language=}, y el estilo visual está configurado globalmente por la plantilla.

\begin{lstlisting}[language=Python, caption={Mi código en Python}, label={code:mifuncion}]
def saludo(nombre):
    print(f"Hola, {nombre}!")
\end{lstlisting}

También es posible insertar código en línea con \lstinline`\lstinline|...|` (sustituyendo \lstinline{|} por cualquier carácter delimitador que no aparezca en el código):

\subsection{Glosario y acrónimos}

El paquete \textit{glossaries} permite gestionar un glosario de términos y una lista de acrónimos. Los términos se definen en el preámbulo y se referencian en el texto con \lstinline|\gls{}|.

La primera vez que aparece un acrónimo en el texto, se muestra su forma expandida seguida de la abreviatura entre paréntesis: \gls{tcp} o \gls{udp}. Las apariciones posteriores muestran únicamente la abreviatura: \gls{tcp} o \gls{udp}. Para forzar siempre la forma larga puede usar \lstinline|\glsfirst{}|, y para la forma corta \lstinline|\acrshort{}|.

\subsection{Referencias cruzadas y citas}

\subsubsection*{Referencias cruzadas}

Cualquier elemento etiquetado con \lstinline|\label{etiqueta}| puede referenciarse con \lstinline|\ref{etiqueta}|. Por ejemplo:

\begin{itemize}\itemsep 0em
    \item La ecuación de Einstein: Ecuación~\ref{eq:einstein}.
    \item Las ecuaciones de Maxwell: Ecuaciones~\ref{eq:gauss} y~\ref{eq:ampere}.
    \item La tabla de protocolos: Tabla~\ref{tab:protocolos}.
    \item La figura del logo: Figura~\ref{fig:logoetsi-cat}.
\end{itemize}

El símbolo \lstinline|~| inserta un espacio irrompible (no divide el texto en dos líneas) entre la palabra y el número, lo que es buena práctica tipográfica.

\subsubsection*{Citas bibliográficas}

Las referencias bibliográficas se gestionan con \textsc{Bib}\TeX{}. Defina sus entradas en el fichero \textit{bibliografia.bib} y puede citarlas en el texto con \lstinline|\cite{clave}|. La plantilla usa el estilo \lstinline{IEEEtran}, habitual en ingeniería \cite{latexproject}.